
\documentclass{article}
\usepackage{colortbl}
\usepackage{makecell}
\usepackage{multirow}
\usepackage{supertabular}

\begin{document}

\newcounter{utterance}

\centering \large Interaction Transcript for game `codenames', experiment `opponent\_difficult', episode 2 with reflexion{-}llama{-}3.1{-}8b{-}instruct.
\vspace{24pt}

{ \footnotesize  \setcounter{utterance}{1}
\setlength{\tabcolsep}{0pt}
\begin{supertabular}{c@{$\;$}|p{.15\linewidth}@{}p{.15\linewidth}p{.15\linewidth}p{.15\linewidth}p{.15\linewidth}p{.15\linewidth}}
    \# & $\;$A & \multicolumn{4}{c}{Game Master} & $\;\:$B\\
    \hline

    \theutterance \stepcounter{utterance}  
    & & & \multicolumn{2}{p{0.3\linewidth}}{
        \cellcolor[rgb]{0.9,0.9,0.9}{
            \makecell[{{p{\linewidth}}}]{
                \texttt{\tiny{[GM$|$GM]}}
                {'hidden': {'team': ['pass', 'spine', 'bottle', 'berry'], 'innocent': ['hospital', 'fall', 'embassy', 'fire', 'wave'], 'opponent': [], 'assassin': ['belt']}, 'revealed': {'team': {'team': [], 'innocent': [], 'opponent': [], 'assassin': []}, 'opponent': {'team': [], 'innocent': [], 'opponent': [], 'assassin': []}}}
            }
        }
    }
    & & \\ \\

    \theutterance \stepcounter{utterance}  
    & & \multicolumn{4}{p{0.6\linewidth}}{
        \cellcolor[rgb]{0.9,0.9,0.9}{
            \makecell[{{p{\linewidth}}}]{
                \texttt{\tiny{[P1$\langle$GM]}}
                \texttt{You have attempted this same task before and did not succeed. These are the plans you wrote after those attempts.} \\
\texttt{Trial \#1: New plan: } \\
\\ 
\texttt{In my previous attempts, I focused on providing clues that were semantically related to multiple words, but I didn't consider the potential for the opposing team to guess an 'opponent' or 'distractor' word, which would end my teammate's turn. This led to a situation where I was providing clues that were too broad, making it difficult for my teammate to guess all of our team's words before the opposing team had uncovered all of theirs.} \\
\\ 
\texttt{Next time, I will focus on providing clues that are more specific and targeted towards a smaller set of words, while still being semantically related to them. I will also pay closer attention to the words that have been removed from the lists and try to avoid providing clues that could lead to the opposing team guessing an 'opponent' or 'distractor' word.} \\
\\ 
\texttt{For example, instead of providing a clue like "CLUE: animal" that relates to multiple words, I could provide a clue like "CLUE: insect" that is more specific and targeted towards the word "bug". This would give my teammate a better chance of guessing the word correctly and avoid ending their turn prematurely.} \\
\\ 
\texttt{Trial \#2: New plan: } \\
\\ 
\texttt{In my previous attempts, I focused on providing clues that were semantically related to multiple words, but I didn't consider the potential for the opposing team to guess an 'opponent' or 'distractor' word, which would end my teammate's turn. This led to a situation where I was providing clues that were too broad, making it difficult for my teammate to guess all of our team's words before the opposing team had uncovered all of theirs.} \\
\\ 
\texttt{Next time, I will focus on providing clues that are more specific and targeted towards a smaller set of words, while still being semantically related to them. I will also pay closer attention to the words that have been removed from the lists and try to avoid providing clues that could lead to the opposing team guessing an 'opponent' or 'distractor' word.} \\
\\ 
\texttt{To achieve this, I will start by identifying the most distinctive features of each word on my team's list and look for a clue that is closely related to those features. For example, if I notice that the words "spine" and "bottle" both have a strong association with the concept of a "container", I could provide a clue like "CLUE: vessel" that is more specific and targeted towards those words.} \\
\\ 
\texttt{I will also try to avoid providing clues that are too general or could be related to multiple words on the opposing team's list. For example, if I notice that the words "hospital" and "embassy" both have a strong association with the concept of a "building", I would avoid providing a clue like "CLUE: building" that could lead to the opposing team guessing an 'opponent' or 'distractor' word.} \\
\\ 
\texttt{By taking a more targeted and specific approach to providing clues, I hope to improve my chances of winning the game and helping my teammate guess all of our team's words before the opposing team has uncovered all of theirs.} \\
\\ 
\\ 
\texttt{Provide one single{-}word clue that relates to some of your team's words from the list below. You can choose to how many and to which words your clue relates to.} \\
\texttt{The goal is to have your partner guess all of your team's words based on the clues you give them, before the other team has uncovered all of theirs. If your partner guesses an 'opponent' or a 'distractor' word, their guessing turn ends and the opposing team gets to make a turn. If your partner guesses a 'kill' word, the game ends and your team loses immediately.} \\
\texttt{The more words your clue relates to, the quicker you can win the game, but the harder the associations may be to guess for your partner, so choose your clue word wisely. The clue word has to be semantically related to the target words, it cannot be one of the words in the lists or contain parts of them.} \\
\\ 
\texttt{Always give your single{-}word clue and your comma{-}separated list of related target words in the following format and make your answers as short as possible, never include any other text than is required in this form:} \\
\\ 
\texttt{CLUE: <WORD>} \\
\texttt{TARGETS: <TARGETS>} \\
\\ 
\texttt{Your team words are: pass, spine, bottle, berry.} \\
\texttt{Your opponent's team words are: .} \\
\texttt{Distractor words are: hospital, fall, embassy, fire, wave.} \\
\texttt{The kill words are: belt.} \\
            }
        }
    }
    & \\ \\

    \theutterance \stepcounter{utterance}  
    & \multicolumn{4}{p{0.6\linewidth}}{
        \cellcolor[rgb]{0.9,0.9,0.9}{
            \makecell[{{p{\linewidth}}}]{
                \texttt{\tiny{[P1$\rangle$GM]}}
                \texttt{CLUE: container} \\
\texttt{TARGETS: bottle, spine} \\
            }
        }
    }
    & & \\ \\

    \theutterance \stepcounter{utterance}  
    & & & \multicolumn{2}{p{0.3\linewidth}}{
        \cellcolor[rgb]{0.9,0.9,0.9}{
            \makecell[{{p{\linewidth}}}]{
                \texttt{\tiny{[GM$|$GM]}}
                \texttt{container} \\
            }
        }
    }
    & & \\ \\

    \theutterance \stepcounter{utterance}  
    & & & \multicolumn{2}{p{0.3\linewidth}}{
        \cellcolor[rgb]{0.9,0.9,0.9}{
            \makecell[{{p{\linewidth}}}]{
                \texttt{\tiny{[GM$|$GM]}}
                ['bottle', 'spine']
            }
        }
    }
    & & \\ \\

    \theutterance \stepcounter{utterance}  
    & & & \multicolumn{2}{p{0.3\linewidth}}{
        \cellcolor[rgb]{0.9,0.9,0.9}{
            \makecell[{{p{\linewidth}}}]{
                \texttt{\tiny{[GM$|$GM]}}
                {'word': 'bottle', 'assignment': 'team'}
            }
        }
    }
    & & \\ \\

    \theutterance \stepcounter{utterance}  
    & & & \multicolumn{2}{p{0.3\linewidth}}{
        \cellcolor[rgb]{0.9,0.9,0.9}{
            \makecell[{{p{\linewidth}}}]{
                \texttt{\tiny{[GM$|$GM]}}
                {'word': 'spine', 'assignment': 'team'}
            }
        }
    }
    & & \\ \\

    \theutterance \stepcounter{utterance}  
    & & & \multicolumn{2}{p{0.3\linewidth}}{
        \cellcolor[rgb]{0.9,0.9,0.9}{
            \makecell[{{p{\linewidth}}}]{
                \texttt{\tiny{[GM$|$GM]}}
                \texttt{opponent has won} \\
            }
        }
    }
    & & \\ \\

    \theutterance \stepcounter{utterance}  
    & \multicolumn{4}{p{0.6\linewidth}}{
        \cellcolor[rgb]{0.95,0.95,0.95}{
            \makecell[{{p{\linewidth}}}]{
                \texttt{\tiny{[P2$\rightarrow$P1]}}
                \texttt{Post{-}episode Reflexion (not sent to GM):} \\
\texttt{New plan: } \\
\\ 
\texttt{In my previous attempts, I focused on providing clues that were semantically related to multiple words, but I didn't consider the potential for the opposing team to guess an 'opponent' or 'distractor' word, which would end my teammate's turn. This led to a situation where I was providing clues that were too broad, making it difficult for my teammate to guess all of our team's words before the opposing team had uncovered all of theirs.} \\
\\ 
\texttt{Next time, I will focus on providing clues that are more specific and targeted towards a smaller set of words, while still being semantically related to them. I will also pay closer attention to the words that have been removed from the lists and try to avoid providing clues that could lead to the opposing team guessing an 'opponent' or 'distractor' word.} \\
\\ 
\texttt{To achieve this, I will start by identifying the most distinctive features of each word on my team's list and look for a clue that is closely related to those features. For example, if I notice that the words "spine" and "bottle" both have a strong association with the concept of a "container", I could provide a clue like "CLUE: vessel" that is more specific and targeted towards those words.} \\
\\ 
\texttt{I will also try to avoid providing clues that are too general or could be related to multiple words on the opposing team's list. For example, if I notice that the words "hospital" and "embassy" both have a strong association with the concept of a "building", I would avoid providing a clue like "CLUE: building" that could lead to the opposing team guessing an 'opponent' or 'distractor' word.} \\
\\ 
\texttt{Additionally, I will pay closer attention to the words that have been removed from the lists and try to avoid providing clues that could lead to the opposing team guessing a 'kill' word. For example, if I notice that the word "belt" has been removed from the list, I would avoid providing a clue like "CLUE: strap" that could lead to the opposing team guessing the 'kill' word.} \\
\\ 
\texttt{By taking a more targeted and specific approach to providing clues, I hope to improve my chances of winning the game and helping my teammate guess all of our team's words before the opposing team has uncovered all of theirs.} \\
            }
        }
    }
    & & \\ \\

\end{supertabular}
}

\end{document}
